\subsection{Architecture, Setup and Adaptations}\label{sec:model}
The model follows Fin-GAN, with the architecture and training setup taken from the reference implementation of
\textcite{vuleticFinGANForecastingClassifying2024}. The adaptations required for the TTF futures dataset are described in
Sec.~\ref{sec:adaptations}.

\subsubsection{Model Architecture}
The ForGAN architecture described in Sec.~\ref{sec:forgan} was used without structural modifications. In $G$, the condition passed through an LSTM layer, described in Sec.~\ref{sec:rnn}. The final hidden state was concatenated with the noise vector $z$. The result passed through a dense layer of width $R_G+N$, then a ReLU activation and another dense layer of width one, where $R_G$ is the LSTM hidden size and $N$ is the noise dimension.
In $D$, the candidate value ($\hat{x}_{t+1}$ or $x_{t+1}$) was appended to the condition, which gave a sequence of size $l+1$. The result passed through an LSTM layer with the hidden size $R_D$. The final hidden state passed through a dense layer with a sigmoid activation (Sec.~\ref{sec:nn}). The network definitions are in \texttt{FinGAN.py}.


Following \textcite{vuleticFinGANForecastingClassifying2024}, $N=8$, $R_G=8$, $R_D=8$ and $l=10$. One value was forecast per condition,
$\hat{x}_{t+1}$. However, the parameters $N$, $R_G$, $R_D$, and $l$ were varied in Sec.~\ref{sec:configurations}. The normal Xavier initialization described in Sec.~\ref{sec:fin-gan} was used.

\subsubsection{Training configuration}\label{sec:training-config}
The RMSProp optimizer, described in Sec.~\ref{sec:nn-training}, was used for both $G$ and $D$, with a learning rate of 0.0001. One discriminator update was performed for each generator update, as described in Sec.~\ref{sec:basic-gan}. The weights were updated after each minibatch of size $n_{\text{batch}}=100$. The training data was shuffled at the start of each epoch.

The input to $G$ and $D$ was standardized using the mean and standard deviation of the first 100 rows of the training window matrix. Using the earliest rows avoids look-ahead bias. As described in Sec.~\ref{sec:data-split}, the training split starts in January 2010. This means that the reference batch covers a period of relatively low volatility. The scaling parameters are not representative of the whole dataset, especially the 2021 to 2022 period of high volatility. The effect of this choice was tested separately, as described in Sec.~\ref{sec:configurations}.

The training was performed in two phases. In the first phase, both $G$ and $D$ were trained on cross-entropy alone for $n_{\text{grad}}$ epochs. The gradient norms were recorded and $\alpha$, $\beta$, $\gamma$ and $\delta$ were computed as described in Sec.~\ref{sec:fin-gan}. In the second phase, the networks were trained for $n_{\text{epochs}}$ further epochs, separately for each of the nine cost function combinations listed in Sec.~\ref{sec:configurations}. An additional tenth branch was trained on cross-entropy alone for $n_{\text{epochs}}$ further epochs, giving the ForGAN baseline used for comparison in Chapter~\ref{sec:exp}. The effective training length was $n_{\text{grad}} + n_{\text{epochs}}$, since the network weights were updated in both phases. The values $n_{\text{grad}}$ and $n_{\text{epochs}}$ were varied across configurations, as described in Sec.~\ref{sec:configurations}.

\textcite{vuleticFinGANForecastingClassifying2024} select the cost function combination and the corresponding $G$ with the highest $\mathrm{SR}_{w}$ on the validation data split (Sec.~\ref{sec:fin-gan}). $\mathrm{SR}_{w}$ measures directional accuracy, while option prices depend on the shape of the return distribution (Sec.~\ref{sec:options}). Therefore, the configuration and the cost function branch were selected on the distributional statistics of Sec.~\ref{sec:evaluation} and their spread across seeds on the test split. The training procedure is implemented in \texttt{run\_ablation.py}.

\subsubsection{Adaptations}\label{sec:adaptations}
The implementation is based on the reference Fin-GAN code published by \textcite{vuleticFinGANForecastingClassifying2024}. The network definitions, the gradient norm matching procedure and the cost function terms of Sec.~\ref{sec:fin-gan} were adopted without
modification. The training loop for the branch trained on cross-entropy alone (Sec.~\ref{sec:training-config}) does not call the backward propagation on the generator loss, so no gradient reaches the generator. The training loops for the other branches do call it. To ensure that the ForGAN baseline is trained under the same procedure, the backward propagation on the generator loss was added to the corresponding training loop.

The data structure with intraday and overnight returns was not adopted. As described in Sec.~\ref{sec:data-pipeline}, the dataset was constructed from daily log returns of the front-month settlement prices of TTF futures. The data loading functions were modified accordingly. Additionally, a separate evaluation module was implemented to compute $\mathrm{SR}_{w}$, the forecast accuracy and the distributional statistics, as described in Sec.~\ref{sec:evaluation}. A new $\mathrm{SR}_{w}$ function was required, because the original implementation sums two $\mathrm{wPnL}^{i}_{2}$ values per day, as described in Sec.~\ref{sec:fin-gan}. This is not applicable to the daily TTF data.

$G$ produces one value per noise sample, $\hat{x}_{t+1}$. To generate price paths, recursive rollout was implemented, as described in Sec.~\ref{sec:pricing-setup}.
